\begin{textsample}{NEG Dim 2 – 2010\_c5 – Score -2.00 – 2010\_c5/enade\_2010\_odonto\_d1b7c708cc.txt}  \label{ex:f2_neg_105}
CAPÍTULO 5 DISTRIBUIÇÃO DOS CONCEITOS 5.1 PANORAMA NACIONAL DA DISTRIBUIÇÃO DOS CONCEITOS A Tabela 5.1 apresenta a quantidade e distribuição de cursos de Odontologia participantes do ENADE/2010 por faixa de conceito e Grande Região.
Observando -se os dados da Tabela 5.1, nota -se que, dos 189 cursos participantes, cerca de 37,1 \% ( 70 cursos ) classificaram -se com conceito 3, o valor modal.
Este conceito modal foi identificado também em quatro das cinco regiões, excetuando -se a região Norte, onde o conceito modal foi 2 ( 30,0 \% ).
Nas demais regiões, os percentuais que variaram de 35,7 \%, na região Centro-Oeste a 40,0 \%, na Sul.
Em termos nacionais, o conceito 4 foi o segundo mais frequente ( 24,3 \% correspondendo a 46 cursos ), e o conceito 2, o terceiro ( 20,1 \% correspondendo a 38 cursos ).
Houve, ainda, nove cursos ( 4,8 \% ) que receberam o conceito 5 e cinco cursos ( 2,6 \% ) que receberam o conceito 1.
Dos 189 cursos, 21 ( 11,1 \% ) ficaram sem conceito ( SC ).
A região Norte participou com 20 cursos, correspondentes a 10,6 \% do total nacional.
Nesta região, o conceito 2 foi atribuído a seis cursos dos 20 participantes, equivalentes a 30,0 \% do total regional.
Outros cinco cursos ( 25,0 \% ) ficaram sem conceito, enquanto três cursos receberam o conceito 1, quatro o conceito 3 e dois o conceito 4.
A região Nordeste participou com 33 cursos, 17,5 \% do total nacional.
Destes, 13 cursos, 39,3 \% em termos regionais, obtiveram conceito 3.
Os conceitos 2 e 4 foram atribuídos a, respectivamente, seis e nove cursos ( 18,2 \% e 27,3 \% ).
O conceito 5 foi atribuído a dois cursos ( 6,1 \% ).
Nessa região, três cursos ( 9,1 \% ) ficaram sem conceito.
Na região Sudeste observou -se a predominância do conceito 3 ( 39,1 \% ), correspondentes a 34 dos 87 cursos participantes.
O conceito 2 foi atribuído a 19 cursos ( 21,8 \% ) e o conceito 4, a 20 ( 23,0 \% ).
O conceito 1 foi atribuído a dois cursos ( 2,3 \% ) e o conceito 5, a quatro ( 4,6 \% ).
Do total de cursos da região, oito ( 9,2 \% ) ficaram sem conceito.
Dos 35 cursos participantes da região Sul, 14 ( 40,0 \% ) obtiveram conceito 3, o conceito modal.
O conceito 2 foi atribuído a quatro cursos ( 11,4 \% ) e o conceito 4, a dez ( 28,6 \% ).
Outros três cursos ( 8,6 \% ) receberam conceito 5.
Os quatro cursos restantes ficaram sem conceito.
Na região Centro-Oeste foram 14 os cursos participantes, dos quais um ficou sem conceito.
O conceito modal foi \textbf{distribuído} pelos conceitos 3 e 4, ambos com cinco cursos ( 35,7 \% para cada ).
Três cursos obtiveram conceito 2 e um deles não recebeu conceito. 5.2 CONCEITOS POR CATEGORIA ADMINISTRATIVA E POR GRANDE REGIÃO A Tabela 5.2 apresenta a distribuição dos cursos participantes do ENADE/2010 de Odontologia, por Categoria Administrativa, de acordo com os conceitos por eles alcançados, segundo as Grandes Regiões brasileiras.
Dos 189 cursos participantes, 54 ( 28,6 \% ) eram ministrados em instituições públicas e 135 ( 71,4 \% ) em privadas.
De acordo com as informações da Tabela 5.2, as instituições públicas foram as que apresentaram, em termos nacionais, um quantitativo maior de cursos com conceito 5, com cinco cursos neste tipo de IES em nove que receberam este conceito.
Dos 54 cursos participantes de IES públicas, o conceito 4 foi atribuído a 32 cursos, o conceito modal.
Classificaram -se, ainda, dois cursos com conceito 2, cinco com conceito 5, doze com conceito 3 e um com conceito 1.
Ficaram sem conceito dois cursos em IES públicas da Área.
Na rede privada, o conceito modal foi 3, com 58 cursos.
Entre os demais cursos participantes, 36 receberam conceito 2 e 14, conceito 4.
O conceito 1 foi atribuído a quatro cursos e o conceito 5, também a quatro.
Nesta categoria, 19 cursos ficaram sem conceito.
Na análise por região, observa -se que, na região Norte, as instituições privadas participaram com 16 cursos ( 80 \% do total da região ), dos quais cinco ficaram sem conceito.
O conceito modal para as instituições privadas na região Norte foi 2 ( com seis cursos ).
Os demais foram avaliados com conceito 1 ( dois cursos ), conceito 3 ( dois cursos ) e conceito 4 ( um curso ).
As instituições públicas participaram com quatro cursos.
O conceito modal para esta Categoria Administrativa, na região Norte, foi 3 ( com dois cursos ).
Os demais foram avaliados com conceito 1 ( um curso ) e conceito 4 ( um curso ).
Na região Nordeste, a rede privada concentrou 16 dos 33 cursos participantes, 48,5 \% do total da região.
Nesta Categoria Administrativa na região, \textbf{predominaram} os cursos com conceito 2 ( oito cursos ).
Os demais foram avaliados com os conceitos 1 ( cinco cursos ), 3 ( um curso ) e 5 ( dois deles ).
As instituições públicas nesta região participaram com 17 cursos.
O conceito modal para esta Categoria Administrativa foi 3 ( com oito cursos ).
Os demais foram avaliados com os conceitos 1 ( um curso ), 2, ( cinco cursos ), 4 ( dois cursos ) e 5 ( um curso ).
Na região Sudeste, a proporção de cursos da rede privada, 79,3 \%, foi quase tão elevada quanto na região Norte, correspondendo a 69 dos 87 cursos participantes.
Nesta categoria na região Sudeste, o conceito modal foi 3 ( 30 cursos ).
Os demais foram avaliados com conceito 1 ( dois cursos ), conceito 2 ( 18 cursos ), conceito 4 ( sete cursos ) e conceito 5 ( quatro cursos ).
Na combinação de Categoria Administrativa e Grande Região, oito cursos não receberam conceito.
Entre os 18 cursos em instituições públicas no Sudeste, a categoria modal foi a que recebeu conceito 4 ( 13 cursos ).
Os demais foram avaliados com os conceitos 2 ( um curso ) e 3 ( quatro cursos ).
As instituições privadas concentraram 23 dos 35 cursos participantes da região Sul, cerca de 66,0 \% do total regional.
Nesse tipo de Categoria Administrativa, 13 cursos ficaram classificados no conceito modal, 3.
Os demais foram avaliados com conceito 2 ( quatro cursos ) e conceito 4 ( três cursos ).
Na combinação de Categoria Administrativa e Grande Região, três cursos não receberam conceito.
As instituições públicas participaram com 12 cursos, dos quais sete receberam o conceito modal, 4.
Os demais foram avaliados com conceito 3 ( um curso ) e conceito 5 ( três cursos ), sendo que um curso ficou sem conceito.
Na região Centro-Oeste, 11 dos 14 cursos participantes eram de instituições privadas ( 78,6 \% em termos regionais ).
Destes, cinco concentraram -se no conceito modal, 3 e um ficou sem conceito.
Os demais receberam conceito 2 ( três cursos ) e conceito 4 ( dois cursos ).
Na rede pública, os três cursos na região Centro-Oeste foram avaliados com conceito 4, o conceito modal. 5.3 CONCEITOS POR ORGANIZAÇÃO ACADÊMICA E POR GRANDE REGIÃO Na Tabela 5.3 encontra -se a distribuição dos conceitos atribuídos aos cursos participantes do ENADE/2010 na Área de Odontologia, por Organização Acadêmica, segundo as Grandes Regiões brasileiras.
Dos 189 cursos participantes, 120 eram oferecidos em Universidades, 25 em Centros Universitários e 44 em Faculdades.
De acordo com os dados apresentados, a maioria dos cursos avaliados com conceito 5 eram vinculados a Universidades ( oito dos nove cursos com este conceito ).
Esse tipo de Organização Acadêmica teve conceito 4 como modal ( com 46 cursos ).
Os demais cursos foram avaliados com conceito 1 ( um curso ), conceito 2 ( 17 cursos ), conceito 3 ( 42 cursos ) e conceito 5 ( oito cursos, como já mencionado ), além de seis cursos que ficaram sem conceito.
Entre os 25 cursos em Centros Universitários, o conceito modal também foi 3 ( com 12 cursos ), e três cursos ficaram sem conceito.
Os outros cursos neste tipo de Organização Acadêmica receberam os conceitos 1 ( dois cursos ), 2 ( sete cursos ) e 5 ( um curso ).
Nas Faculdades, dos 44 cursos, 12 ficaram sem conceito e 16 receberam o conceito 3, o valor modal.
Os outros cursos neste tipo de Organização Acadêmica receberam conceito 1 ( dois cursos ) e conceito 2 ( 14 cursos ).
Considerando -se separadamente as regiões brasileiras, verifica -se que, na região Norte, as Universidades concentraram quatro dos 20 cursos participantes.
Dos cursos em Universidades, dois se classificaram com conceito 3 e os outros dois, com conceito 4.
Os Centros Universitários da região Norte foram representados, também, por quatro cursos, sendo dois avaliados com conceito 2, o valor modal, e os demais com os conceitos 1 e 3.
Na região Norte, as Faculdades concentraram o maior número de cursos, 12, sendo que cinco ficaram sem conceito.
O conceito 1 foi atribuído a dois cursos, o conceito 2 a quatro cursos, sendo o conceito modal, e o conceito 3 a um curso.
Na região Nordeste, as Universidades participaram com 20 dos 33 cursos na Área de Odontologia da região.
Entre os cursos em Universidades, um ficou sem conceito.
O conceito modal foi 4 ( nove cursos ).
Os demais cursos receberam conceito 2 ( um curso ), conceito 3 ( sete cursos ) e conceito 5 ( dois cursos ).
Os Centros Universitários contaram com três cursos participantes na região Nordeste.
Apenas um curso não recebeu conceito e os outros dois ficaram entre os conceitos 2 e 3 ( nesse caso, o conceito modal ).
As Faculdades na região Nordeste foram representadas por dez cursos, dos quais um não recebeu conceito.
O conceito modal foi 3 ( com cinco cursos ).
Os quatro cursos restantes foram classificados com conceito 2.
Na região Sudeste, as Universidades concentraram 59 dos 87 cursos da região.
Entre os cursos em Universidades, o conceito modal foi 4 ( com 20 cursos ) e três ficaram sem conceito.
Os demais cursos receberam conceito 1 ( um curso ), conceito 2 ( 14 cursos ), conceito 3 ( 18 cursos ) e conceito 5 ( três cursos ).
Os Centros Universitários participaram com 13 cursos na região Sudeste, dos quais seis obtiveram o conceito modal, 3, e dois ficaram sem conceito.
Os demais receberam conceito 1 ( um curso ), conceito 2 ( três cursos ), e conceito 5 ( um curso ).
As Faculdades foram representadas por 15 cursos na região Sudeste, que se \textbf{distribuíram} em conceito 2 ( dois cursos ), e o conceito modal, 3, com dez cursos.
Três deles ficaram sem conceito.
Dos 35 cursos da região Sul, 29 eram de Universidades, para os quais o conceito modal foi 3, com 12 cursos.
Nesse tipo de organização, dois cursos ficaram sem conceito e os demais receberam conceito 2 ( dois cursos ), conceito 4 ( dez cursos ) e conceito 5 ( três cursos ).
Os Centros Universitários da região Sul tiveram participação de dois cursos e estes foram avaliados com conceito 3.
Foram quatro os cursos vinculados a Faculdades na região Sul, e destes, dois ficaram sem conceito.
Os dois restantes receberam conceito 2.
Na região Centro-Oeste, oito dos 14 cursos eram de Universidades.
Nesse tipo de organização, o conceito modal foi 4 ( com cinco cursos ).
Os três cursos remanescentes foram avaliados com conceito 3.
Os Centros Universitários da região Centro-Oeste contaram com três cursos, dois receberam conceito 3, o valor modal e um recebeu conceito 2.
Dos três cursos em Faculdades da região Centro-Oeste, um ficou sem conceito e dois receberam o conceito modal, 2.

% matched lemmas: distribuir, predominar
\end{textsample}
